\section{\texorpdfstring{$G$}{G}-injective matrix product tensors}%
\label{sec:peps_mps_g_injective}

A matrix product tensor is a family of operators $A^i$ on a bond space
carrying the representation $U$. Its virtual system is the space of
operators $X$ on the bond space, on which $G$ acts by conjugation
$X\mapsto U_gXU_g^{-1}$, and identifying the virtual vector
$\sum x_{\alpha\beta}\ket{\alpha,\beta}$ with the operator
$X=\sum x_{\alpha\beta}\ket{\beta}\bra{\alpha}$ gives
$\mathcal{P}(A)X=\sum_i\tr[A^iX]\ket{i}$
\cite[Section~``MPS and PEPS'']{Schuch2010PEPS}.

\begin{definition}[$G$-injective matrix product tensor]%
    \label{def:peps_mps_g_injective}
    \lean{TNLean.PEPS.mpsSiteMap, TNLean.PEPS.mpsSiteMap_apply}
    \leanok
    \uses{def:peps_g_injective}
    A matrix product tensor $A$ is $G$-injective if the map
    $\mathcal{P}(A)X=\sum_i\tr[A^iX]\ket{i}$ is $G$-injective for the
    conjugation action of $U$ on the operators on the bond space.
\end{definition}

\begin{lemma}[$G$-injective matrix product tensors]%
    \label{thm:peps_mps_g_injective_source_form}
    \lean{TNLean.PEPS.mpsSiteMap_comp_linHom,
        TNLean.PEPS.comp_smulRight_dual,
        TNLean.PEPS.eq_of_forall_trace_mul_eq,
        TNLean.PEPS.mpsSiteMap_comp_linHom_eq_iff,
        TNLean.PEPS.isGInjective_mpsSiteMap_iff}
    \leanok
    \uses{def:peps_mps_g_injective}
    A matrix product tensor $A$ is $G$-injective if and only if
    \begin{enumerate}
        \item $U_gA^iU_g^{-1}=A^i$ for all $i$ and $g$, and
        \item $\mathcal{P}(A)$ has a left inverse on the commutant
              $\mathcal{S}=\{X\mid [X,U_g]=0\}$: there is a linear map $L$
              with $L\,\mathcal{P}(A)X=X$ for every $X\in\mathcal{S}$.
    \end{enumerate}
    For a unitary representation $U_g^{-1}=U_g^\dagger$, and this is the
    source's definition
    \cite[Section~``MPS: the $G$-injective case'']{Schuch2010PEPS}.
\end{lemma}

\begin{proof}\leanok
    \uses{thm:peps_group_averaging}
    Since $\tr[A^iU_gXU_g^{-1}]=\tr[U_g^{-1}A^iU_gX]$, invariance of
    $\mathcal{P}(A)$ under conjugation by $U_g$ says that $U_g^{-1}A^iU_g$
    and $A^i$ have the same trace pairing with every operator $X$; pairing
    with the rank-one operators $\ket{v}\bra{f}$ shows that they are equal.
    The operators fixed by conjugation are the commutant, and a linear map
    is injective on a subspace if and only if it has a left inverse there.
\end{proof}

\begin{lemma}[Stability under concatenation]%
    \label{thm:peps_mps_g_injective_concatenation}
    \lean{TNLean.PEPS.concatTensor, TNLean.PEPS.concatTensor_apply,
        TNLean.PEPS.concatLeftInverse,
        TNLean.PEPS.apply_mpsSiteMap,
        TNLean.PEPS.averageMap_linHom_apply_mul,
        TNLean.PEPS.sum_leftInverse_mul_mul,
        TNLean.PEPS.concatLeftInverse_comp_mpsSiteMap,
        TNLean.PEPS.IsGInjective.mpsSiteMap_concatTensor}
    \leanok
    \uses{def:peps_mps_g_injective, def:peps_delta_operator}
    Let $A$ and $B$ be $G$-injective matrix product tensors for the same
    representation, with left inverses $L_A$ and $L_B$ satisfying
    $L_A\mathcal{P}(A)=L_B\mathcal{P}(B)=\sigma$. Then $C^{ij}=A^iB^j$ is
    $G$-injective, and the map
    \begin{align*}
        L_C\ket{i,j}=|G|\,\big(L_B\ket{j}\big)\,\Delta\,\big(L_A\ket{i}\big)
    \end{align*}
    satisfies $L_C\mathcal{P}(C)=\sigma$
    \cite[Section~``MPS: the $G$-injective case'']{Schuch2010PEPS}.
\end{lemma}

\begin{proof}\leanok
    \uses{thm:peps_mps_g_injective_source_form,
        thm:peps_g_injective_left_inverse, thm:peps_delta_expansion}
    Invariance of $C$ follows from
    $U_gA^iB^jU_g^{-1}=(U_gA^iU_g^{-1})(U_gB^jU_g^{-1})$. For the left
    inverse, first $\sum_i\tr[A^iY]\,L_A\ket{i}=\sigma(Y)$ for every $Y$.
    Second, for every operator $M$,
    \begin{align*}
        \sum_j\big(L_B\ket{j}\big)\,M\,B^j
        =\frac{1}{|G|}\sum_h\tr[U_{h^{-1}}M]\,U_h ;
    \end{align*}
    writing $M$ in a basis $\{b_k\}$ with coordinate functionals $b^*_k$,
    the left-hand side applied to $w$ is
    $\sum_k\sigma\big(\ket{w}\bra{b^*_k}\big)Mb_k$, which is the right-hand
    side applied to $w$. Taking $Y=B^jX$ in the first identity and
    $M=\Delta U_g$ in the second,
    \begin{align*}
        L_C\mathcal{P}(C)X
        =|G|\sum_j\big(L_B\ket{j}\big)\,\Delta\,\sigma(B^jX)
        =\sum_g\Big(\sum_j\big(L_B\ket{j}\big)\,\Delta U_g\,B^j\Big)XU_g^{-1}
        =\frac{1}{|G|}\sum_gU_gXU_g^{-1},
    \end{align*}
    the last step by the expansion
    $\sum_h\tr[U_{h^{-1}}\Delta U_g]U_h=U_g$.
\end{proof}

\begin{lemma}[Blocking a $G$-injective tensor]%
    \label{cor:peps_mps_g_injective_blocking}
    \lean{TNLean.PEPS.IsGInjective.mpsSiteMap_reindex,
        TNLean.PEPS.wordTensor,
        TNLean.PEPS.IsGInjective.mpsSiteMap_wordTensor,
        MPSTensor.toLin'_evalWord,
        MPSTensor.isGInjective_mpsSiteMap_blockTensor}
    \leanok
    \uses{def:peps_mps_g_injective, def:blocked_tensor}
    If $A$ is $G$-injective, then the tensor obtained by blocking $L\ge1$
    sites, $(i_1,\dots,i_L)\mapsto A^{i_1}\cdots A^{i_L}$, is $G$-injective.
\end{lemma}

\begin{proof}\leanok
    \uses{thm:peps_mps_g_injective_concatenation}
    Induction on $L$: the blocked tensor on $L+1$ sites is the concatenation
    of $A$ with the blocked tensor on $L$ sites, up to relabelling the
    physical index, and relabelling preserves $G$-injectivity.
\end{proof}

For $G$-injective tensors the parent-Hamiltonian argument runs as in the
injective case, with the left inverses replaced by left inverses up to the
twirl $\sigma$
\cite[Section~``Parent Hamiltonians'']{Schuch2010PEPS}. In the next two
theorems $A$ and $B$ are $G$-injective for the same representation $U$; three-leg
tensors $M$, $N$ carry one physical and two virtual legs, and the operators $X$,
$M$, $N$, $U_g$ sit on a bond.

\begin{lemma}[Twirling the closing operator]%
    \label{lem:peps_mps_g_injective_twirl}
    \lean{TNLean.PEPS.mpsSiteMap_averageMap,
        TNLean.PEPS.averageMap_mul_of_conj_eq,
        TNLean.PEPS.mul_comm_of_conj_eq,
        TNLean.PEPS.concatTensor_conj_eq}
    \leanok
    \uses{def:peps_mps_g_injective}
    If $U_gC^iU_g^{-1}=C^i$ for all $i$ and $g$, then
    $\tr[C^i\sigma(X)]=\tr[C^iX]$ for every operator $X$: an operator closing
    an invariant tensor may be replaced by its twirl without changing the
    state. In addition, $\sigma(YZ)=\sigma(Y)Z$ whenever $Z$ commutes with every
    $U_g$, and the concatenation of two invariant tensors is invariant
    \cite[Section~``Parent Hamiltonians'']{Schuch2010PEPS}.
\end{lemma}

\begin{proof}\leanok
    \uses{thm:peps_mps_g_injective_source_form, thm:peps_group_averaging}
    Invariance of $C$ is invariance of the map $X\mapsto(\tr[C^iX])_i$ under
    conjugation, and an invariant map does not see the average over the
    group. For the second claim, $U_gYZU_g^{-1}=(U_gYU_g^{-1})Z$.
\end{proof}

\begin{lemma}[The concatenated left inverse on traces]%
    \label{lem:peps_mps_g_injective_left_inverse_traces}
    \lean{TNLean.PEPS.concatLeftInverse_apply_trace_mul}
    \leanok
    \uses{thm:peps_mps_g_injective_concatenation}
    Let $L_A$ be any linear map and let $L_B\mathcal{P}(B)=\sigma$. The map
    $L_C\ket{j,k}=|G|\,(L_A\ket{k})\,\Delta\,(L_B\ket{j})$ of the concatenation
    $C^{jk}=B^jA^k$ sends the vector $(\tr[B^jY_k])_{j,k}$ to
    $|G|\sum_k(L_A\ket{k})\,\Delta\,\sigma(Y_k)$, for every family of
    operators $Y_k$.
\end{lemma}

\begin{proof}\leanok
    Summing over $j$ first, $\sum_j\tr[B^jY_k]\,L_B\ket{j}
        =L_B\mathcal{P}(B)Y_k=\sigma(Y_k)$.
\end{proof}

\begin{lemma}[Intersection property for $G$-injective tensors]%
    \label{thm:peps_mps_g_injective_intersection}
    \lean{TNLean.PEPS.IsGInjective.intersection_property}
    \leanok
    \uses{thm:peps_mps_g_injective_source_form,
        thm:peps_mps_g_injective_concatenation}
    \begin{align*}
        \Big\{\sum_{ijk}\tr[A^iB^jM^k]\ket{ijk}\ \Big|\ M\Big\}\cap
        \Big\{\sum_{ijk}\tr[N^iB^jA^k]\ket{ijk}\ \Big|\ N\Big\}
        =\Big\{\sum_{ijk}\tr[A^iB^jA^kX]\ket{ijk}\ \Big|\ X\Big\}
    \end{align*}
    \cite[Section~``Parent Hamiltonians'']{Schuch2010PEPS}.
\end{lemma}

\begin{proof}\leanok
    \uses{lem:peps_mps_g_injective_twirl,
        lem:peps_mps_g_injective_left_inverse_traces}
    The right-hand side lies in both sets, with $M^k=A^kX$ and $N^i=XA^i$.
    Conversely, let $\tr[A^iB^jM^k]=\tr[N^iB^jA^k]$ for all $i,j,k$. Fix
    left inverses with $L_A\mathcal{P}(A)=L_B\mathcal{P}(B)=\sigma$, and let
    $L_C$ be the left inverse of $C^{jk}=B^jA^k$, so that
    $L_C\mathcal{P}(C)=\sigma$. For fixed $i$ the vector
    $(\tr[B^jA^kN^i])_{j,k}$ equals $(\tr[B^jM^kA^i])_{j,k}$, and applying
    $L_C$ gives
    \begin{align*}
        \sigma(N^i)=|G|\sum_k(L_A\ket{k})\,\Delta\,\sigma(M^kA^i)=XA^i,
        \qquad X=|G|\sum_k(L_A\ket{k})\,\Delta\,\sigma(M^k),
    \end{align*}
    because $A^i$ commutes with every $U_g$. Since $C$ is invariant,
    $\tr[N^iB^jA^k]=\tr[B^jA^k\sigma(N^i)]=\tr[A^iB^jA^kX]$.
\end{proof}

\begin{lemma}[Closure property for $G$-injective tensors]%
    \label{thm:peps_mps_g_injective_closure}
    \lean{TNLean.PEPS.IsGInjective.eq_sum_of_closure,
        TNLean.PEPS.IsGInjective.closure_property}
    \leanok
    \uses{thm:peps_mps_g_injective_source_form,
        thm:peps_mps_g_injective_concatenation, def:peps_delta_operator}
    \begin{align*}
        \Big\{\sum_{ij}\tr[B^iA^jM]\ket{ij}\ \Big|\ M\Big\}\cap
        \Big\{\sum_{ij}\tr[B^iNA^j]\ket{ij}\ \Big|\ N\Big\}
        =\Big\{\sum_g\lambda_g\sum_{ij}\tr[B^iA^jU_g]\ket{ij}\ \Big|\
        \lambda_g\in\C\Big\}.
    \end{align*}
    A vector in the intersection, given by $M$ and $N$, has
    $\lambda_g=\tr[U_{g^{-1}}\sigma(N)\Delta]$
    \cite[Section~``Parent Hamiltonians'']{Schuch2010PEPS}.
\end{lemma}

\begin{proof}\leanok
    \uses{lem:peps_mps_g_injective_twirl,
        lem:peps_mps_g_injective_left_inverse_traces}
    The right-hand side lies in both sets, with $M=N=\sum_g\lambda_gU_g$,
    since $U_g$ commutes with $A^j$. Conversely, let
    $\tr[B^iA^jM]=\tr[B^iNA^j]$ for all $i,j$, and let $L_C$ be the left
    inverse of $C^{ij}=B^iA^j$. Applying $L_C$ to the vector
    $(\tr[B^iNA^j])_{i,j}$, and then the identity
    $\sum_j(L_A\ket{j})\,Y\,A^j=\frac{1}{|G|}\sum_h\tr[U_{h^{-1}}Y]\,U_h$ with
    $Y=\Delta\sigma(N)$,
    \begin{align*}
        \sigma(M)=|G|\sum_j(L_A\ket{j})\,\Delta\,\sigma(N)\,A^j
        =\sum_h\tr[U_{h^{-1}}\Delta\sigma(N)]\,U_h .
    \end{align*}
    Since $C$ is invariant, $\tr[B^iA^jM]=\tr[B^iA^j\sigma(M)]$, and
    $\tr[U_{h^{-1}}\Delta\sigma(N)]=\tr[U_{h^{-1}}\sigma(N)\Delta]$ because
    $\sigma(N)$ commutes with $U_{h^{-1}}$.
\end{proof}

The two properties give the parent Hamiltonian of a $G$-injective tensor, as
the intersection and closure properties of injective tensors give Theorem 3.5
of the source \cite[Section~``Parent Hamiltonians'']{Schuch2010PEPS}. For a
tensor $A$ and a ring of $N$ sites, write $\mathcal G_{N,L}(A)$ for the space
of Definition~\ref{def:chain_ground_space}, the common kernel of the parent
interactions on all windows of $L$ consecutive sites of the ring, and write
$\ket{\mathcal M(A|K)}=\sum_{\sigma}\tr[A^{\sigma_1}\cdots A^{\sigma_N}K]
    \ket{\sigma}$ for the matrix product state with closure $K$. The source's
two-site interaction is $L=2$. In the rest of this section $A$ is
$G$-injective for the representation $U$.

\begin{lemma}[The open chain of a $G$-injective tensor]%
    \label{lem:peps_mps_g_injective_open_chain}
    \lean{MPSTensor.mem_groundSpace_add_two_of_isGInjective,
        MPSTensor.chainGroundSpace_le_groundSpace_of_isGInjective}
    \leanok
    \uses{thm:peps_mps_g_injective_intersection,
        cor:peps_mps_g_injective_blocking, def:chain_ground_space}
    Let $K\ge1$. A vector on $K+2$ sites whose restrictions to the first and to
    the last $K+1$ sites lie in
    $\{\sum_{\sigma}\tr[A^{\sigma_1}\cdots A^{\sigma_{K+1}}X]\ket{\sigma}\mid X\}$
    lies in the corresponding space on $K+2$ sites. Consequently, if
    $2\le L\le N$, every vector of $\mathcal G_{N,L}(A)$ has the form
    $\sum_{\sigma}\tr[A^{\sigma_1}\cdots A^{\sigma_N}X]\ket{\sigma}$
    \cite[Section~``Parent Hamiltonians'']{Schuch2010PEPS}.
\end{lemma}

\begin{proof}\leanok
    Fixing the last letter $k$ writes the vector as $\tr[A^iB^wM^k]$ and fixing
    the first letter $i$ as $\tr[N^iB^wA^k]$, where $B$ is the tensor of $K$
    blocked sites, which is $G$-injective by
    Lemma~\ref{cor:peps_mps_g_injective_blocking}. The intersection
    property gives the form $\tr[A^iB^wA^kX]$. Iterating along the
    non-wrapping windows of the ring gives the second claim, as in the proof of
    Theorem 3.5 of the source.
\end{proof}

\begin{lemma}[Parent Hamiltonians of $G$-injective tensors]%
    \label{thm:peps_mps_g_injective_parent_hamiltonian}
    \lean{MPSTensor.groundSpaceMap_mem_span_closure_of_isGInjective,
        MPSTensor.groundSpaceMap_rep_mem_chainGroundSpace,
        MPSTensor.chainGroundSpace_eq_span_closure_of_isGInjective}
    \leanok
    \uses{lem:peps_mps_g_injective_open_chain,
        thm:peps_mps_g_injective_closure}
    If $2\le L\le N$, then
    \begin{align*}
        \mathcal G_{N,L}(A)=\spn_{\C}\bigl\{\ket{\mathcal M(A|U_g)}\ \big|\
        g\in G\bigr\}.
    \end{align*}
    In particular, for the two-site interaction on a ring of $N\ge2$ sites
    the frustration-free ground space of the parent Hamiltonian is spanned by
    the states $\ket{\mathcal M(A|U_g)}$
    \cite[Theorem~4.11]{Schuch2010PEPS}.
\end{lemma}

\begin{proof}\leanok
    Each $U_g$ commutes with every $A^i$, so it can be moved to any bond, and
    every window of $\ket{\mathcal M(A|U_g)}$ lies in the local ground space.
    Conversely, by Lemma~\ref{lem:peps_mps_g_injective_open_chain} a vector
    of $\mathcal G_{N,L}(A)$ is $\tr[A^{\sigma_1}\cdots A^{\sigma_N}X]$, and
    its translate by one site is $\tr[A^{\sigma_1}\cdots A^{\sigma_N}Y]$. Hence
    $\tr[A^aB^wX]=\tr[A^aYB^w]$, where $B$ is the tensor of the last $N-1$
    blocked sites. The closure property, applied to $A$ and $B$, writes the
    vector as a combination of the states $\ket{\mathcal M(A|U_g)}$.
\end{proof}

\begin{lemma}[Twirled representation operators]%
    \label{lem:peps_twirled_representation}
    \lean{Representation.rep_mul_charProjector,
        Representation.apply_conj_of_mem_irreducibleCharacters,
        Representation.sum_rep_conj_apply_of_mem,
        Representation.averageMap_linHom_rep,
        Representation.linearIndependent_charProjector,
        Representation.linearIndependent_of_isSemiRegular,
        Representation.linearIndependent_averageMap_conjClasses}
    \leanok
    \uses{thm:peps_isotypic_projectors, thm:peps_group_averaging,
        thm:peps_semiregular_trace}
    The isotypic projectors $P_i$ commute with every $U_g$, are linearly
    independent, and
    \begin{align*}
        \sigma(U_g)=\frac{1}{|G|}\sum_hU_{hgh^{-1}}
        =\sum_{i=1}^{r}\frac{\chi_i(g)}{\chi_i(1)}\,P_i .
    \end{align*}
    If $U$ is semi-regular, then the operators $U_g$ are linearly independent,
    and so are the twirls $\sigma(U_g)$ of one representative $g$ of each
    conjugacy class \cite[Section~``Parent Hamiltonians'']{Schuch2010PEPS}.
\end{lemma}

\begin{proof}\leanok
    On an irreducible subrepresentation with character $\chi_i$ the operator
    $\sum_hU_{hgh^{-1}}$ commutes with the representation, so by Schur's lemma
    it is a scalar, and its trace $|G|\chi_i(g)$ fixes the scalar as
    $|G|\chi_i(g)/\chi_i(1)$; the projector $P_i$ acts there as the identity
    and every other projector as zero. Evaluating $\sum_i\mu_iP_i$ on a nonzero
    vector of such a subrepresentation gives $\mu_i$ times the vector. For
    semi-regular $U$, the functional $X\mapsto\tr[U_{g^{-1}}X\Delta]$ takes the
    value $\delta_{g,h}$ on $U_h$ by
    Lemma~\ref{thm:peps_semiregular_trace}; on $\sigma(U_k)$ it is
    $\frac{1}{|G|}$ times the number of $h$ with $hkh^{-1}=g$, which vanishes
    unless $k$ is conjugate to $g$ and is positive for $k=g$.
\end{proof}

\begin{theorem}[Structure of the ground space of a $G$-injective tensor]%
    \label{thm:peps_mps_g_injective_ground_space_structure}
    \lean{MPSTensor.groundSpaceMap_toMatrix'_averageMap,
        MPSTensor.chainGroundSpace_eq_span_charProjector_of_isGInjective,
        MPSTensor.linearIndependent_groundSpaceMap_charProjector_of_isGInjective,
        MPSTensor.finrank_chainGroundSpace_of_isGInjective,
        MPSTensor.chainGroundSpace_eq_span_conjClasses_of_isGInjective,
        MPSTensor.linearIndependent_groundSpaceMap_conjClasses_of_isSemiRegular,
        MPSTensor.finrank_chainGroundSpace_of_isSemiRegular,
        MPSTensor.card_irreducibleCharacterFinset_of_isSemiRegular}
    \leanok
    \uses{thm:peps_mps_g_injective_parent_hamiltonian,
        lem:peps_twirled_representation, def:peps_semiregular}
    Let $2\le L\le N$, and let $\chi_1,\dots,\chi_I$ be the distinct
    characters of the irreducible representations occurring in $U$, with
    isotypic projectors $P_i$. The states $\ket{\mathcal M(A|P_i)}$ are
    linearly independent and span $\mathcal G_{N,L}(A)$, so this space has
    dimension $I$. If $U$ is semi-regular, then $I$ is the number of conjugacy
    classes of $G$, and for any choice of one representative $g$ of each
    conjugacy class the states $\ket{\mathcal M(A|U_g)}$ are linearly
    independent and span $\mathcal G_{N,L}(A)$
    \cite[Theorem~4.12]{Schuch2010PEPS}.
\end{theorem}

\begin{proof}\leanok
    Since $A$ is invariant, $\ket{\mathcal M(A|K)}=\ket{\mathcal M(A|\sigma(K))}$
    for every closure $K$; in particular the closures of one conjugacy class
    give the same state. By Lemma~\ref{lem:peps_twirled_representation} each
    $\ket{\mathcal M(A|U_g)}$ is a combination of the states
    $\ket{\mathcal M(A|P_i)}$, and conversely each $P_i$ is a combination of
    the $U_g$, so Lemma~\ref{thm:peps_mps_g_injective_parent_hamiltonian}
    gives the spanning claims. The tensor of $N$ blocked sites is
    $G$-injective, so $K\mapsto\ket{\mathcal M(A|K)}$ is injective on the
    operators commuting with every $U_g$. These contain the projectors $P_i$
    and the twirls $\sigma(U_g)$, and
    Lemma~\ref{lem:peps_twirled_representation} gives the linear independence
    of both families. Computing the dimension of $\mathcal G_{2,2}(A)$ in the
    two ways gives $I$ equal to the number of conjugacy classes. The source
    instead quotes that the numbers of conjugacy classes and of irreducible
    representations agree, and passes from the $U_g$ to the $P_i$ through the
    completeness of the characters in the class functions.
\end{proof}
